% SGA 3, Expose VI_B, Loop 1 English reconstruction.
% Scope: Definition 6.6, Remark 6.7, Example 6.8, and paragraph
% 6.9 in full, including source editor's note 67.
% Authority: Polo--Gille Exp6B-13oct24.pdf, component-local page
% 49 (combined-reader page 390; re-edition header 379;
% original-margin printed page 362).
% Section 7 begins on local page 50 and is excluded.
% OCR and Reinhold English Markdown are comparison witnesses only.
% This unit contains no diagrams.

\medskip
\phantomsection\label{def:sga3-VIB-6.6}%
\noindent\textbf{定义 6.6.---}
设 \(G\) 为群值 \(S\)-函子, \(H\) 为 \(G\) 的群
子 \(S\)-函子. 若
\[
  \underline{\operatorname{Norm}}_G H=G
\]
(相应地, 若
\(\underline{\operatorname{Centr}}_G H=G\); 相应地, 若对
每个 \(S\)-概形 \(T\) 及每个自同构
\[
  a\in\operatorname{Aut}_{T\text{-}\mathrm{gr}}(G_T)
\]
有 \(a(H_T)\subset H_T\)), 则称 \(H\) 在 \(G\) 中\emph{不变}
(相应地, 为\emph{中心的}; 相应地, 为\emph{特征的}). 换言之, 对每个
\(S\)-概形 \(T\), \(G(T)\) 的子群 \(H(T)\) 在 \(G(T)\) 中正规
(相应地, 在 \(G(T)\) 中为中心的; 相应地,
在 \(G_T\) 的每个自同构下不变).

\noindent\emph{注意.}
若 \(H\) 为中心的 (相应地, 特征的), 则它
不变.

\medskip
\phantomsection\label{rem:sga3-VIB-6.7}%
\noindent\textbf{注 6.7.---}
设 \(G\) 和 \(H\) 为两个 \(S\)-群, \(u:H\to G\) 为
单态射. 群 \(H\) 在 \(G\) 中不变 (相应地,
为中心的) 当且仅当态射
\[
  \mu\circ
  \bigl(c\circ\operatorname{pr}_2,\,
        \mu\circ(u\times\operatorname{id}_G)\bigr)
  :H\times_S G\longrightarrow G
\]
(其定义为: 对 \(g\in G(S')\) 和 \(h\in H(S')\),
\[
  (h,g)\longmapsto g^{-1}hg
\]
) 经由 \(u\) 分解
(相应地, 等于 \(u\circ\operatorname{pr}_1\)).
群 \(H\) 在 \(G\) 中为特征的当且仅当, 对
每个 \(S\)-概形 \(T\) 及 \(G_T\) 的每个 \(T\)-群自同构 \(a\),
态射 \(a\circ u(T)\) 经由 \(u(T)\) 分解.

\medskip
\phantomsection\label{ex:sga3-VIB-6.8}%
\noindent\textbf{例 6.8.---}
设 \(G\) 为群值 \(S\)-函子. 则
\(\underline{\operatorname{Centr}}G\) 为特征的且
为中心的. 若对每个 \(s\in S\), \(G_s\) 可表, 则
\(G^0\) 为特征的. 这由定义和
\hyperref[prop:sga3-VIB-3.3]{3.3} 得到.

\medskip
\phantomsection\label{par:sga3-VIB-6.9}%
\noindent\textbf{6.9.---}%
\footnote{\label{note:sga3-VIB-67}编者注: 本小节
为新增内容.}
在 [RG71], I, \SGARefLink{cor:sga3-VIA-3.3.3}{3.3.3} 中, 作者引入了
\emph{纯 \(S\)-概形}这一几何概念, 它在 \(S\) 上关于
平展拓扑是局部的; 精确定义见\emph{同上}.
这里只记录以下几点:
\begin{enumerate}
\item[\textup{a)}]\phantomsection\label{point:sga3-VIB-6.9-a}
  (\emph{同上}, 定理 3.3.5) 若 \(B\) 为平坦的
  有限呈示 \(A\)-代数, 则 \(B\) 为投射
  \(A\)-模当且仅当
  \(\operatorname{Spec}(B)\) 在
  \(\operatorname{Spec}(A)\) 上纯.
\item[\textup{b)}]\phantomsection\label{point:sga3-VIB-6.9-b}
  因此, 若 \(X\to S\) 局部有限
  呈示、平坦且纯, 则 \(X\) 在
  \(S\) 上本质自由.
\item[\textup{c)}]\phantomsection\label{point:sga3-VIB-6.9-c}
  (\emph{同上}, 3.3.4\textup{(iii)}) 若 \(X\to S\)
  局部有限型、平坦、纤维几何不可约
  且无嵌入分支, 则 \(X\) 在
  \(S\) 上纯.
\end{enumerate}
域上的每个局部有限型群概形都是
Cohen--Macaulay 的
(参见 \hyperref[prop:sga3-VIA-1.1.1]{VI\(_A\), 1.1.1}),
因而无嵌入分支. 特别地得到:
\begin{enumerate}
\item[\textup{d)}]\phantomsection\label{point:sga3-VIB-6.9-d}
  每个局部有限呈示、平坦且纤维连通的
  \(S\)-群概形 \(G\) 都在
  \(S\) 上纯, 从而在 \(S\) 上本质自由.
\end{enumerate}
因此可以结合上面的 \textup{(b)} 和 \textup{(d)} 的结果,
使用 \hyperref[ex:sga3-VIB-6.2.4-e]{6.2.4\textup{(e)}} 的
所有陈述.

% 翻译统计: 源行数 106; 输出行数 106 (含本统计行); 行内/行间数学片段 67; 翻译占位 0. 底本疑点另报主线程.
